% =====================================================================
%  CAPITULO 11 - CIRCUITOS DE PRIMEIRA E SEGUNDA ORDEM
% =====================================================================
\section{Circuitos de Primeira e Segunda Ordem}

\begin{conceito}[Mapa de resolução]
Este capítulo trata do comportamento transitório de circuitos com elementos
armazenadores de energia. A solução deve partir da \textbf{leitura do circuito},
e não apenas da substituição de valores em fórmulas.

\textbf{Roteiro para circuitos com chaveamento:}
\begin{enumerate}[leftmargin=1.8em,itemsep=2pt,topsep=3pt]
\item desenhe ou interprete o circuito para $t<0$ e obtenha o regime permanente;
\item use a continuidade para determinar $v_C(0^+)$ e/ou $i_L(0^+)$;
\item redesenhe o circuito para $t>0$, determine o valor final e a resistência
      vista pelo elemento armazenador;
\item somente então escreva a resposta transitória e verifique unidades,
      condições iniciais e comportamento para $t\to\infty$.
\end{enumerate}

\textbf{Primeira ordem:}
\[
x(t)=x(\infty)+\big[x(0^+)-x(\infty)\big]e^{-t/\tau}.
\]
Para circuitos RC, $\tau=R_{\mathrm{eq}}C$ e a tensão do capacitor é contínua:
$v_C(0^+)=v_C(0^-)$. Para circuitos RL, $\tau=L/R_{\mathrm{eq}}$ e a corrente do
indutor é contínua: $i_L(0^+)=i_L(0^-)$. Para obter $R_{\mathrm{eq}}$, observe a
resistência vista pelo elemento armazenador com as fontes independentes
suprimidas; fontes dependentes, quando existirem, permanecem ativas.

\textbf{Segunda ordem:} dois elementos independentes de armazenamento podem
produzir uma resposta de segunda ordem. Eles não precisam ser necessariamente um
capacitor e um indutor: duas capacitâncias ou duas indutâncias também podem gerar
dois estados quando não podem ser reduzidas a um único elemento equivalente.
Antes de resolver a equação de segunda ordem, determine as condições iniciais,
os valores finais e, quando necessário, $\dot v(0^+)$ ou $\dot i(0^+)$ pelas leis
de Kirchhoff.

\textbf{RLC em série:}
\[
\alpha=\frac{R}{2L},\qquad
\omega_0=\frac{1}{\sqrt{LC}},\qquad
\zeta=\frac{\alpha}{\omega_0}.
\]
Se $\alpha<\omega_0$, a resposta é subamortecida e
$\omega_d=\sqrt{\omega_0^2-\alpha^2}$. Se $\alpha=\omega_0$, há amortecimento
crítico. Se $\alpha>\omega_0$, a resposta é superamortecida. No RLC em série,
$R_{\mathrm{crit}}=2\sqrt{L/C}$. No RLC em paralelo, para a tensão natural,
$\alpha=1/(2RC)$.
\end{conceito}

\legendaniveis

% Os bancos são separados por nível para manter o capítulo organizado.
% N1 permanece conceitual; N2--N4 recebem uma trilha adicional com muitos
% circuitos, chaveamentos, formas de onda, projeto e diagnóstico.
% =====================================================================
% CAP. 11 - NIVEL 1
% =====================================================================
\blocofixacao

\begin{questao}[1]{}
Um capacitor ideal possui $v_C(0^-)=\SI{7}{\volt}$. No instante $t=0$ ocorre uma
comutação no circuito. Qual é o valor de $v_C(0^+)$?
\resposta{$v_C(0^+)=\SI{7}{\volt}$}
\resolucao{A tensão em um capacitor ideal não pode variar instantaneamente, pois
$i_C=C\,\mathrm{d}v_C/\mathrm{d}t$ exigiria corrente impulsiva infinita. Logo,
$v_C(0^+)=v_C(0^-)=\SI{7}{\volt}$.}
\end{questao}

\begin{questao}[1]{}
Um indutor ideal conduz $i_L(0^-)=\SI{350}{\milli\ampere}$. Uma chave muda de
posição em $t=0$. Determine $i_L(0^+)$.
\resposta{$i_L(0^+)=\SI{350}{\milli\ampere}$}
\resolucao{A corrente de um indutor ideal é contínua. Como
$v_L=L\,\mathrm{d}i_L/\mathrm{d}t$, uma mudança instantânea de corrente exigiria
tensão impulsiva infinita. Portanto $i_L(0^+)=i_L(0^-)$.}
\end{questao}

\begin{questao}[1]{}
Um circuito RC tem $R=\SI{2}{\kilo\ohm}$ e $C=\SI{50}{\micro\farad}$. Determine
a constante de tempo.
\resposta{$\tau=\SI{100}{\milli\second}$}
\resolucao{$\tau=RC=(\pot{2}{3})(\pot{50}{-6})=\SI{0.1}{\second}
=\SI{100}{\milli\second}$.}
\end{questao}

\begin{questao}[1]{}
Um circuito RL possui $L=\SI{80}{\milli\henry}$ e resistência equivalente
$R_{\mathrm{eq}}=\SI{20}{\ohm}$. Calcule $\tau$.
\resposta{$\tau=\SI{4}{\milli\second}$}
\resolucao{$\tau=L/R_{\mathrm{eq}}=0{,}08/20=\SI{0.004}{\second}
=\SI{4}{\milli\second}$.}
\end{questao}

\begin{questao}[1]{}
Um capacitor inicialmente descarregado é carregado por uma fonte CC através de
um resistor. Após um intervalo igual a uma constante de tempo, aproximadamente
qual fração da variação total de tensão já ocorreu?
\begin{alternativas}[3]
\item \SI{36.8}{\percent}
\item \SI{50}{\percent}
\item \SI{63.2}{\percent}
\item \SI{86.5}{\percent}
\item \SI{99.3}{\percent}
\end{alternativas}
\resposta{(c)}
\resolucao{Para uma carga partindo de zero,
$v_C(\tau)=V_f(1-e^{-1})\approx0{,}632V_f$. Portanto cerca de
\SI{63.2}{\percent} da variação total já ocorreu.}
\end{questao}

\begin{questao}[1]{}
Uma resposta natural de primeira ordem é proporcional a $e^{-t/\tau}$. Qual
percentual do valor inicial resta após $3\tau$?
\resposta{$e^{-3}\times100\%\approx\SI{4.98}{\percent}$}
\resolucao{$e^{-3}\approx0{,}0498$. Assim resta cerca de \SI{4.98}{\percent} do
valor inicial.}
\end{questao}

\begin{questao}[1]{}
Ao calcular a resistência equivalente vista por um capacitor para obter a
constante de tempo, como devem ser tratadas as fontes independentes ideais?
\begin{alternativas}
\item fontes de tensão abertas e fontes de corrente em curto.
\item todas as fontes abertas.
\item fontes de tensão em curto e fontes de corrente abertas.
\item todas as fontes em curto.
\item as fontes não devem ser modificadas.
\end{alternativas}
\resposta{(c)}
\resolucao{Para calcular $R_{\mathrm{eq}}$ com fontes independentes suprimidas,
uma fonte ideal de tensão é substituída por curto-circuito e uma fonte ideal de
corrente por circuito aberto. Fontes dependentes, quando existirem, permanecem
ativas.}
\end{questao}

\begin{questao}[1]{}
Um circuito ideal possui um capacitor e um indutor independentes entre si como
únicos elementos armazenadores de energia. Qual é, em geral, a ordem máxima do
circuito?
\resposta{Segunda ordem}
\resolucao{Cada variável independente de armazenamento de energia pode introduzir
um estado. Um capacitor e um indutor independentes produzem duas variáveis de
estado, caracterizando um circuito de segunda ordem.}
\end{questao}

\begin{questao}[1]{}
Para $L=\SI{0.2}{\henry}$ e $C=\SI{50}{\micro\farad}$, determine a frequência
natural não amortecida $\omega_0$.
\resposta{$\omega_0\approx\SI{316.2}{\radian\per\second}$}
\resolucao{$LC=0{,}2\cdot\pot{50}{-6}=\pot{1}{-5}$ e
$\omega_0=1/\sqrt{LC}=1/\sqrt{\pot{1}{-5}}\approx\SI{316.2}{\radian\per\second}$.}
\end{questao}

\begin{questao}[1]{}
Em um RLC em série, $R=\SI{40}{\ohm}$, $L=\SI{0.1}{\henry}$ e
$C=\SI{100}{\micro\farad}$. Classifique o amortecimento.
\resposta{Subamortecido}
\resolucao{$\alpha=R/(2L)=40/0{,}2=\SI{200}{\per\second}$ e
$\omega_0=1/\sqrt{0{,}1\cdot\pot{100}{-6}}\approx\SI{316.2}{\radian\per\second}$.
Como $\alpha<\omega_0$, a resposta é subamortecida.}
\end{questao}

\begin{questao}[1]{}
Determine a resistência crítica de um RLC em série com $L=\SI{50}{\milli\henry}$ e
$C=\SI{200}{\micro\farad}$.
\resposta{$R_{\mathrm{crit}}\approx\SI{31.6}{\ohm}$}
\resolucao{$R_{\mathrm{crit}}=2\sqrt{L/C}=2\sqrt{0{,}05/(\pot{200}{-6})}
=2\sqrt{250}\approx\SI{31.6}{\ohm}$.}
\end{questao}

\gabarito[Nível 1]

% =====================================================================
% CAP. 11 - NIVEL 2
% =====================================================================
\blocoaplicacao

\begin{questao}[2]{}
A tensão de um capacitor muda de $\SI{8}{\volt}$ para um valor final de
$\SI{2}{\volt}$ através de $R=\SI{4}{\kilo\ohm}$ e
$C=\SI{25}{\micro\farad}$. Determine $v_C(\SI{0.2}{\second})$.
\resposta{$v_C(\SI{0.2}{\second})\approx\SI{2.81}{\volt}$}
\resolucao{$\tau=RC=\SI{0.1}{\second}$. Logo
\[
v_C(t)=2+(8-2)e^{-t/0{,}1}.
\]
Em $t=\SI{0.2}{\second}$, $v_C=2+6e^{-2}\approx\SI{2.81}{\volt}$.}
\end{questao}

\begin{questao}[2]{}
A corrente de um indutor passa de $\SI{1}{\ampere}$ para um valor final de
$\SI{4}{\ampere}$. Sabendo que $L=\SI{0.6}{\henry}$ e
$R_{\mathrm{eq}}=\SI{30}{\ohm}$, calcule $i_L(\SI{40}{\milli\second})$.
\resposta{$i_L\approx\SI{3.59}{\ampere}$}
\resolucao{$\tau=L/R=0{,}6/30=\SI{20}{\milli\second}$. Assim
$i_L(t)=4+(1-4)e^{-t/\tau}$. Para $t=2\tau$,
$i_L=4-3e^{-2}\approx\SI{3.59}{\ampere}$.}
\end{questao}

\begin{questao}[2]{}
Um circuito RC inicialmente descarregado possui $R=\SI{10}{\kilo\ohm}$ e
$C=\SI{10}{\micro\farad}$. Quanto tempo é necessário para atingir
\SI{90}{\percent} do valor final?
\resposta{$t\approx\SI{230}{\milli\second}$}
\resolucao{$\tau=\SI{0.1}{\second}$. De
$0{,}9=1-e^{-t/\tau}$ resulta $e^{-t/\tau}=0{,}1$ e
$t=\tau\ln 10\approx\SI{0.230}{\second}$.}
\end{questao}

\begin{questao}[2]{}
Um capacitor descarrega de $\SI{12}{\volt}$ para $\SI{3}{\volt}$ com
$\tau=\SI{50}{\milli\second}$. Determine o tempo decorrido.
\resposta{$t\approx\SI{69.3}{\milli\second}$}
\resolucao{$3=12e^{-t/\tau}$, então $e^{-t/\tau}=1/4$ e
$t=\tau\ln4=\SI{50}{\milli\second}\cdot1{,}386\approx\SI{69.3}{\milli\second}$.}
\end{questao}

\begin{questao}[2]{}
A corrente natural de um circuito RL cai de $\SI{2}{\ampere}$ para
$\SI{0.2}{\ampere}$. Se $\tau=\SI{25}{\milli\second}$, determine o intervalo de
tempo.
\resposta{$t\approx\SI{57.6}{\milli\second}$}
\resolucao{$0{,}2=2e^{-t/\tau}$ implica $e^{-t/\tau}=0{,}1$. Logo
$t=\tau\ln10\approx25\cdot2{,}303=\SI{57.6}{\milli\second}$.}
\end{questao}

\begin{questao}[2]{}
No circuito abaixo, a chave conecta a fonte em $t=0$ e o capacitor está
inicialmente descarregado. Determine $v_C(\tau)$ e a constante de tempo.

\circuitoq{%
\begin{circuitikz}[scale=1,transform shape,line width=0.8pt]
\draw (0,0) to[battery1,l_={$\SI{20}{\volt}$}] (0,3)
      to[R,l={$R_1=\SI{6}{\kilo\ohm}$}] (3,3) coordinate(a);
\draw (a) to[R,l={$R_2=\SI{3}{\kilo\ohm}$}] (3,0) -- (0,0);
\draw (a) -- (5,3) to[C,l={$C=\SI{50}{\micro\farad}$}] (5,0) -- (3,0);
\end{circuitikz}}

\resposta{$\tau=\SI{100}{\milli\second}$; $v_C(\tau)\approx\SI{4.21}{\volt}$}
\resolucao{Em regime permanente, a tensão no capacitor corresponde à tensão do divisor:
$V_f=20\cdot3/(6+3)=\SI{6.67}{\volt}$. Suprimindo a fonte,
$R_{\mathrm{eq}}=R_1\parallel R_2=\SI{2}{\kilo\ohm}$, então
$\tau=R_{\mathrm{eq}}C=\SI{0.1}{\second}$. Como $v_C(0)=0$,
$v_C(\tau)=V_f(1-e^{-1})\approx\SI{4.21}{\volt}$.}
\end{questao}

\begin{questao}[2]{}
Após uma comutação, um capacitor possui $v_C(0^+)=\SI{5}{\volt}$ e tende a
$-\SI{10}{\volt}$, com $\tau=\SI{20}{\milli\second}$. Determine
$v_C(\SI{10}{\milli\second})$.
\resposta{$v_C\approx-\SI{0.90}{\volt}$}
\resolucao{$v_C(t)=-10+[5-(-10)]e^{-t/\tau}=-10+15e^{-t/\tau}$.
Para $t/\tau=0{,}5$, obtém-se $v_C\approx-10+15e^{-0{,}5}
\approx-\SI{0.90}{\volt}$.}
\end{questao}

\begin{questao}[2]{}
Um indutor de $\SI{40}{\milli\henry}$ está conectado, após uma
comutação, a uma rede cuja resistência equivalente vista em seus terminais é
$R_{\mathrm{eq}}=\SI{4}{\ohm}$. Sua corrente inicial é $\SI{0.5}{\ampere}$ e o valor final é
$\SI{2}{\ampere}$. Calcule a corrente após $\SI{20}{\milli\second}$.
\resposta{$i_L\approx\SI{1.80}{\ampere}$}
\resolucao{$\tau=L/R=\SI{10}{\milli\second}$ e
$i_L(t)=2+(0{,}5-2)e^{-t/\tau}$. Em $t=2\tau$,
$i_L=2-1{,}5e^{-2}\approx\SI{1.80}{\ampere}$.}
\end{questao}

\begin{questao}[2]{}
Um circuito RLC em série tem $R=\SI{10}{\ohm}$, $L=\SI{50}{\milli\henry}$ e
$C=\SI{200}{\micro\farad}$. Determine $\alpha$, $\omega_0$, $\omega_d$ e o tipo
de resposta.
\resposta{$\alpha=\SI{100}{\per\second}$; $\omega_0\approx\SI{316.2}{\radian\per\second}$;
$\omega_d=\SI{300}{\radian\per\second}$; subamortecida}
\resolucao{$\alpha=R/(2L)=10/0{,}1=100$. Além disso,
$\omega_0=1/\sqrt{0{,}05\cdot\pot{200}{-6}}\approx316{,}2$. Como
$\alpha<\omega_0$, a resposta é subamortecida e
$\omega_d=\sqrt{316{,}2^2-100^2}=\SI{300}{\radian\per\second}$.}
\end{questao}

\begin{questao}[2]{}
Para um circuito RLC em série com $R=\SI{40}{\ohm}$, $L=\SI{50}{\milli\henry}$ e
$C=\SI{200}{\micro\farad}$, classifique a resposta natural.
\resposta{Superamortecida}
\resolucao{$\alpha=40/(2\cdot0{,}05)=\SI{400}{\per\second}$, enquanto
$\omega_0\approx\SI{316.2}{\radian\per\second}$. Como $\alpha>\omega_0$, a
resposta é superamortecida.}
\end{questao}

\begin{questao}[2]{}
Calcule $R_{\mathrm{crit}}$ para um circuito RLC em série com $L=\SI{0.1}{\henry}$ e
$C=\SI{100}{\micro\farad}$.
\resposta{$R_{\mathrm{crit}}\approx\SI{63.25}{\ohm}$}
\resolucao{$R_{\mathrm{crit}}=2\sqrt{L/C}=2\sqrt{0{,}1/\pot{100}{-6}}
=2\sqrt{1000}\approx\SI{63.25}{\ohm}$.}
\end{questao}

\begin{questao}[2]{}
Uma resposta subamortecida de tensão tem a forma
\[
v(t)=e^{-100t}\big(10\cos300t+B\sin300t\big).
\]
Sabendo que $\mathrm{d}v/\mathrm{d}t|_{0^+}=0$, determine $B$.
\resposta{$B\approx\SI{3.33}{\volt}$}
\resolucao{Derivando e avaliando em zero,
$\dot v(0)=-100\cdot10+300B=0$. Portanto
$B=1000/300\approx\SI{3.33}{\volt}$.}
\end{questao}

\begin{questao}[2]{}
Em um capacitor de $C=\SI{100}{\micro\farad}$, a corrente em $0^+$ é
$\SI{20}{\milli\ampere}$ entrando no terminal positivo. Determine
$\mathrm{d}v_C/\mathrm{d}t|_{0^+}$.
\resposta{$\mathrm{d}v_C/\mathrm{d}t|_{0^+}=\SI{200}{\volt\per\second}$}
\resolucao{Pela convenção passiva, $i_C=C\dot v_C$. Logo
$\dot v_C=0{,}02/(\pot{100}{-6})=\SI{200}{\volt\per\second}$.}
\end{questao}

\begin{questao}[2]{}
Um circuito RLC em paralelo possui $R=\SI{200}{\ohm}$, $C=\SI{100}{\micro\farad}$ e
$L=\SI{0.1}{\henry}$. Classifique a resposta natural da tensão.
\resposta{Subamortecida}
\resolucao{No circuito RLC em paralelo, $\alpha=1/(2RC)=1/[2\cdot200\cdot\pot{100}{-6}]
=\SI{25}{\per\second}$. Já $\omega_0\approx\SI{316.2}{\radian\per\second}$.
Como $\alpha<\omega_0$, a resposta é subamortecida.}
\end{questao}

\begin{questao}[2]{}
Em uma resposta de primeira ordem, qual erro percentual em relação ao valor final
permanece após $5\tau$?
\resposta{$\approx\SI{0.674}{\percent}$}
\resolucao{O erro normalizado é $e^{-t/\tau}$. Em $5\tau$,
$e^{-5}\approx0{,}00674$, ou \SI{0.674}{\percent}.}
\end{questao}

\begin{questao}[2]{}
Um circuito LC ideal possui um capacitor de $C=\SI{100}{\micro\farad}$
com tensão inicial de $\SI{10}{\volt}$ e um indutor de
$L=\SI{50}{\milli\henry}$ com corrente inicial nula.
Determine a energia inicial e a corrente máxima ideal.
\resposta{$E=\SI{5}{\milli\joule}$; $I_{\max}\approx\SI{0.447}{\ampere}$}
\resolucao{$E=\tfrac12CV^2=\tfrac12\pot{100}{-6}\cdot100=\SI{5}{\milli\joule}$.
Quando toda a energia está no indutor,
$\tfrac12LI_{\max}^2=E$, então
$I_{\max}=V\sqrt{C/L}=10\sqrt{\pot{100}{-6}/0{,}05}\approx\SI{0.447}{\ampere}$.}
\end{questao}

\begin{questao}[2]{}
Para $L=\SI{10}{\milli\henry}$ e $C=\SI{100}{\micro\farad}$, determine o
período natural ideal $T_0$.
\resposta{$T_0\approx\SI{6.28}{\milli\second}$}
\resolucao{$\omega_0=1/\sqrt{LC}=1/\sqrt{\pot{1}{-6}}=\SI{1000}{\radian\per\second}$.
Logo $T_0=2\pi/\omega_0\approx\SI{6.28}{\milli\second}$.}
\end{questao}

\gabarito[Nível 2]

% =====================================================================
% CAP. 11 - NIVEL 2 - QUESTOES GRAFICAS E INTEGRADAS
% Inspiradas no nivel de variedade de livros-texto, com enunciados originais.
% =====================================================================
\blocoaplicacao

\begin{questao}[2]{}
A chave do circuito permaneceu fechada por muito tempo e \textbf{abre} em
$t=0$. Determine $v_C(0^+)$, a constante de tempo para $t>0$ e
$v_C(\SI{0.30}{\second})$.

\circuitoqq{%
\begin{circuitikz}[scale=0.96,transform shape,line width=0.8pt,every node/.append style={font=\scriptsize}]
\draw (0,0) to[american voltage source,l_={$\SI{18}{\volt}$}] (0,3)
      to[switch] (1.9,3)
      to[R,l={$R_1=\SI{3}{\kilo\ohm}$}] (4.5,3) coordinate(a);
\node[anchor=south] at (0.95,3.30) {abre em $t=0$};
\draw (a) to[R,l_={$R_2=\SI{6}{\kilo\ohm}$}] (4.5,0) -- (0,0);
\draw (a) -- (7.0,3)
      to[C,v^>={$v_C$}] (7.0,0) -- (4.5,0);
\node[anchor=east] at (6.75,1.5) {$C=\SI{100}{\micro\farad}$};
\end{circuitikz}}{Circuito RC com mudança de topologia em $t=0$.}

\resposta{$v_C(0^+)=\SI{12}{\volt}$; $\tau=\SI{0.60}{\second}$;
$v_C(0{,}30\,\mathrm{s})\approx\SI{7.28}{\volt}$}
\resolucao{Para $t<0$, o capacitor está em regime permanente e equivale a um
circuito aberto. Assim,
\[
v_C(0^-)=18\frac{6}{3+6}=\SI{12}{\volt}.
\]
Pela continuidade da tensão no capacitor, $v_C(0^+)=\SI{12}{\volt}$. Quando a
chave abre, a fonte e $R_1$ ficam desconectados do capacitor; a descarga ocorre
somente por $R_2$. Logo
\[
\tau=R_2C=(6\times10^3)(100\times10^{-6})=\SI{0.60}{\second}.
\]
Como $v_C(\infty)=0$,
$v_C(t)=12e^{-t/0{,}60}$ e, em $t=\SI{0.30}{\second}$,
$v_C\approx12e^{-0{,}5}=\SI{7.28}{\volt}$.}
\end{questao}

\begin{questao}[2]{}
No circuito abaixo, a chave permaneceu fechada por muito tempo e abre em
$t=0$. Determine $i_L(0^+)$, a constante de tempo da descarga e
$i_L(\SI{10}{\milli\second})$.

\circuitoqq{%
\begin{circuitikz}[scale=0.96,transform shape,line width=0.8pt,every node/.append style={font=\scriptsize}]
\draw (0,0) to[american voltage source,l_={$\SI{24}{\volt}$}] (0,3)
      to[switch] (1.9,3)
      to[R,l={$R_1=\SI{6}{\ohm}$}] (4.5,3) coordinate(a);
\node[anchor=south] at (0.95,3.30) {abre em $t=0$};
\draw (a) to[L,i>^={$i_L$}] (4.5,0) -- (0,0);
\node[anchor=east] at (4.25,1.5) {$L=\SI{120}{\milli\henry}$};
\draw (a) -- (7.0,3) to[R,l_={$R_2=\SI{18}{\ohm}$}] (7.0,0) -- (4.5,0);
\end{circuitikz}}{Circuito RL com caminho de descarga resistivo.}

\resposta{$i_L(0^+)=\SI{4}{\ampere}$; $\tau\approx\SI{6.67}{\milli\second}$;
$i_L(10\,\mathrm{ms})\approx\SI{0.893}{\ampere}$}
\resolucao{Em regime permanente, o indutor equivale a um curto-circuito; portanto
$R_2$ fica sem tensão e
$i_L(0^-)=24/6=\SI{4}{\ampere}$. Pela continuidade da corrente,
$i_L(0^+)=\SI{4}{\ampere}$. Após a abertura da chave, a fonte e $R_1$ ficam
isolados e o indutor descarrega por $R_2$:
\[
\tau=\frac{L}{R_2}=\frac{0{,}12}{18}\approx\SI{6.67}{\milli\second}.
\]
Assim, $i_L(t)=4e^{-t/\tau}$ e
$i_L(10\,\mathrm{ms})\approx4e^{-10/6{,}67}=\SI{0.893}{\ampere}$.}
\end{questao}

\begin{questao}[2]{}
A fonte de corrente já está ligada há muito tempo. Em $t=0$ a chave fecha,
inserindo $R_2$ em paralelo com o capacitor. Determine $v_C(0^+)$,
$v_C(\infty)$, $\tau$ e escreva $v_C(t)$ para $t>0$.

\circuitoqq{%
\begin{circuitikz}[scale=0.94,transform shape,line width=0.8pt,every node/.append style={font=\scriptsize}]
\draw (0,0) to[american current source,l_={$\SI{2}{\milli\ampere}$}] (0,3) -- (2.2,3) coordinate(a);
\draw (a) to[R,l_={$R_1=\SI{4}{\kilo\ohm}$}] (2.2,0) -- (0,0);
\draw (a) -- (4.7,3) to[C,v^>={$v_C$}] (4.7,0) -- (2.2,0);
\node[anchor=east] at (4.45,1.5) {$C=\SI{50}{\micro\farad}$};
\draw (a) -- (6.0,3) to[switch] (7.5,3)
      to[R,l_={$R_2=\SI{12}{\kilo\ohm}$}] (7.5,0) -- (4.7,0);
\node[anchor=south] at (6.75,3.30) {fecha em $t=0$};
\end{circuitikz}}{Mudança da resistência vista pelo capacitor.}

\resposta{$v_C(0^+)=\SI{8}{\volt}$; $v_C(\infty)=\SI{6}{\volt}$;
$\tau=\SI{150}{\milli\second}$;
$v_C(t)=6+2e^{-t/0{,}15}\ \si{\volt}$}
\resolucao{Antes da comutação, o capacitor está aberto e $R_2$ desconectado. Logo
$v_C(0^-)=I_SR_1=(2\,\mathrm{mA})(4\,\mathrm{k\Omega})=\SI{8}{\volt}$, portanto
$v_C(0^+)=\SI{8}{\volt}$. Após o fechamento, em regime permanente o capacitor
continua aberto e
\[
v_C(\infty)=I_S(R_1\parallel R_2)
=2\,\mathrm{mA}\,(3\,\mathrm{k\Omega})=\SI{6}{\volt}.
\]
Para a constante de tempo, a fonte de corrente independente é aberta:
$R_{\mathrm{eq}}=R_1\parallel R_2=\SI{3}{\kilo\ohm}$ e
$\tau=R_{\mathrm{eq}}C=\SI{0.15}{\second}$. Assim,
$v_C(t)=6+(8-6)e^{-t/0{,}15}$.}
\end{questao}

\begin{questao}[2]{}
A chave fecha em $t=0$, aplicando um degrau de $\SI{20}{\volt}$ ao circuito RLC
em série. Sem resolver ainda a resposta completa, determine $\alpha$, $\omega_0$,
$\omega_d$ e classifique o amortecimento.

\circuitoqq{%
\begin{circuitikz}[scale=0.92,transform shape,line width=0.8pt,every node/.append style={font=\scriptsize}]
\draw (0,0) to[american voltage source,l_={$\SI{20}{\volt}$}] (0,3)
      to[switch] (1.8,3)
      to[R,l={$R=\SI{15}{\ohm}$}] (4.0,3)
      to[L,l={$L=\SI{50}{\milli\henry}$}] (6.5,3)
      to[C] (9.0,3) -- (9.0,0) -- (0,0);
\node[anchor=south] at (0.9,3.30) {fecha em $t=0$};
\node[anchor=south] at (7.75,3.30) {$C=\SI{200}{\micro\farad}$};
\end{circuitikz}}{Circuito RLC em série excitado por degrau.}

\resposta{$\alpha=\SI{150}{\per\second}$;
$\omega_0\approx\SI{316.2}{\radian\per\second}$;
$\omega_d\approx\SI{278.4}{\radian\per\second}$; resposta subamortecida}
\resolucao{Para RLC em série,
\[
\alpha=\frac{R}{2L}=\frac{15}{0{,}10}=\SI{150}{\per\second},\qquad
\omega_0=\frac{1}{\sqrt{LC}}
=\frac{1}{\sqrt{0{,}05\cdot200\times10^{-6}}}
\approx\SI{316.2}{\radian\per\second}.
\]
Como $\alpha<\omega_0$, a resposta é subamortecida e
$\omega_d=\sqrt{\omega_0^2-\alpha^2}\approx\SI{278.4}{\radian\per\second}$.}
\end{questao}

\begin{questao}[2]{}
O circuito RLC em paralelo abaixo parte sem energia armazenada e recebe a fonte
$i_s(t)=3u(t)\,\si{\ampere}$. Determine $v(0^+)$, $i_L(0^+)$,
$\mathrm{d}v/\mathrm{d}t|_{0^+}$, $v(\infty)$ e $i_L(\infty)$.

\circuitoqq{%
\begin{circuitikz}[scale=0.96,transform shape,line width=0.8pt,every node/.append style={font=\scriptsize}]
\draw (0,0) to[american current source] (0,3) -- (7.2,3);
\node[anchor=east] at (-0.25,1.5) {$3u(t)\,\si{\ampere}$};
\draw (2.2,3) to[R] (2.2,0);
\node[anchor=east] at (1.95,1.5) {$R=\SI{6}{\ohm}$};
\draw (4.7,3) to[C,v^>={$v$}] (4.7,0);
\node[anchor=east] at (4.45,1.5) {$C=\SI{0.5}{\farad}$};
\draw (7.2,3) to[L,i>^={$i_L$}] (7.2,0) -- (0,0);
\node[anchor=east] at (6.95,1.5) {$L=\SI{2}{\henry}$};
\end{circuitikz}}{RLC em paralelo: condições iniciais, derivada inicial e regime final.}

\resposta{$v(0^+)=0$; $i_L(0^+)=0$;
$\dot v(0^+)=\SI{6}{\volt\per\second}$;
$v(\infty)=0$; $i_L(\infty)=\SI{3}{\ampere}$}
\resolucao{Como o circuito parte sem energia, $v_C(0^-)=0$ e $i_L(0^-)=0$;
logo $v(0^+)=0$ e $i_L(0^+)=0$. Aplicando KCL em $0^+$,
\[
3=\frac{v}{R}+C\dot v+i_L
\quad\Rightarrow\quad
3=0+0{,}5\dot v+0,
\]
portanto $\dot v(0^+)=\SI{6}{\volt\per\second}$. Em regime CC, o capacitor é
aberto e o indutor é curto-circuito; assim $v(\infty)=0$ e toda a corrente da fonte
flui pelo indutor: $i_L(\infty)=\SI{3}{\ampere}$.}
\end{questao}

\gabarito[Nível 2 --- circuitos e chaveamento]

% =====================================================================
% CAP. 11 - NIVEL 3
% =====================================================================
\blocoaprofundamento

\begin{questao}[3]{}
Um capacitor de $\SI{10}{\micro\farad}$, inicialmente descarregado, é ligado a
$\SI{12}{\volt}$ por $R_1=\SI{2}{\kilo\ohm}$ durante
$\SI{50}{\milli\second}$. Em seguida, a fonte é removida e ele passa a
descarregar por $R_2=\SI{5}{\kilo\ohm}$. Determine sua tensão
$\SI{40}{\milli\second}$ após a segunda comutação.
\resposta{$v_C\approx\SI{4.95}{\volt}$}
\resolucao{Na primeira etapa, $\tau_1=\SI{20}{\milli\second}$ e
$v_C(50\,\mathrm{ms})=12(1-e^{-2{,}5})\approx\SI{11.015}{\volt}$. Esse valor é a
condição inicial da segunda etapa. Agora $\tau_2=\SI{50}{\milli\second}$ e
$v_C=11{,}015e^{-40/50}\approx\SI{4.95}{\volt}$.}
\end{questao}

\begin{questao}[3]{}
Um indutor de $\SI{100}{\milli\henry}$ é energizado por uma fonte de
$\SI{20}{\volt}$ através de $R_1=\SI{10}{\ohm}$ durante
$\SI{15}{\milli\second}$. Depois, a fonte é retirada e o indutor descarrega por
$R_2=\SI{25}{\ohm}$. Qual é a corrente após mais $\SI{8}{\milli\second}$?
\resposta{$i_L\approx\SI{0.210}{\ampere}$}
\resolucao{Na energização, $I_f=20/10=\SI{2}{\ampere}$ e
$\tau_1=0{,}1/10=\SI{10}{\milli\second}$. Assim
$i_L(15\,\mathrm{ms})=2(1-e^{-1{,}5})\approx\SI{1.554}{\ampere}$. Na descarga,
$\tau_2=0{,}1/25=\SI{4}{\milli\second}$, portanto após $2\tau_2$:
$i_L=1{,}554e^{-2}\approx\SI{0.210}{\ampere}$.}
\end{questao}

\begin{questao}[3]{}
Após a supressão das fontes independentes, a resistência vista por um
capacitor de $C=\SI{20}{\micro\farad}$ é composta por
$R_3=\SI{1}{\kilo\ohm}$ em série com o paralelo de
$R_1=\SI{6}{\kilo\ohm}$ e $R_2=\SI{3}{\kilo\ohm}$. Determine a
constante de tempo.

\circuitoq{%
\begin{circuitikz}[scale=1,transform shape,line width=0.8pt]
\draw (0,1.5) to[C,l={$C$},o-] (2,1.5) to[R,l={$R_3$}] (4,1.5) coordinate(a);
\draw (a) -- (4,2.6) to[R,l={$R_1$}] (6,2.6) -- (6,0.4);
\draw (a) -- (4,0.4) to[R,l_={$R_2$}] (6,0.4) -- (6,-0.2) node[ground]{};
\end{circuitikz}}

\resposta{$R_{\mathrm{eq}}=\SI{3}{\kilo\ohm}$; $\tau=\SI{60}{\milli\second}$}
\resolucao{$R_1\parallel R_2=(6\cdot3)/(6+3)=\SI{2}{\kilo\ohm}$. Em série com
$R_3$, $R_{\mathrm{eq}}=\SI{3}{\kilo\ohm}$. Logo
$\tau=R_{\mathrm{eq}}C=\pot{3}{3}\pot{20}{-6}=\SI{60}{\milli\second}$.}
\end{questao}

\begin{questao}[3]{}
Após uma comutação, $v_C(0^+)=-\SI{4}{\volt}$, $v_C(\infty)=\SI{8}{\volt}$,
$R_{\mathrm{eq}}=\SI{2.5}{\kilo\ohm}$ e $C=\SI{40}{\micro\farad}$. Em que
instante a tensão cruza zero?
\resposta{$t\approx\SI{40.5}{\milli\second}$}
\resolucao{$\tau=\SI{0.1}{\second}$ e
$v_C=8+(-4-8)e^{-t/0{,}1}=8-12e^{-t/0{,}1}$. No cruzamento,
$e^{-t/0{,}1}=2/3$. Portanto $t=0{,}1\ln(3/2)\approx\SI{40.5}{\milli\second}$.}
\end{questao}

\begin{questao}[3]{}
Um indutor possui $i_L(0^+)=-\SI{1}{\ampere}$, $i_L(\infty)=\SI{3}{\ampere}$ e
$\tau=\SI{25}{\milli\second}$. Determine quando a corrente muda de sinal.
\resposta{$t\approx\SI{7.19}{\milli\second}$}
\resolucao{$i_L=3+(-1-3)e^{-t/\tau}=3-4e^{-t/\tau}$. Para $i_L=0$,
$e^{-t/\tau}=3/4$. Logo $t=\tau\ln(4/3)\approx\SI{7.19}{\milli\second}$.}
\end{questao}

\begin{questao}[3]{}
Um capacitor com $v_C(0)=\SI{2}{\volt}$ é ligado a uma fonte de
$\SI{10}{\volt}$ através de $R=\SI{2}{\kilo\ohm}$ e
$C=\SI{50}{\micro\farad}$. Determine a corrente inicial e a corrente após
$\SI{150}{\milli\second}$.
\resposta{$i(0^+)=\SI{4}{\milli\ampere}$; $i(150\,\mathrm{ms})\approx\SI{0.893}{\milli\ampere}$}
\resolucao{$\tau=RC=\SI{0.1}{\second}$. A corrente é
$i(t)=(10-2)R^{-1}e^{-t/\tau}=\SI{4}{\milli\ampere}e^{-t/\tau}$. Em
$150\,\mathrm{ms}=1{,}5\tau$, $i\approx4e^{-1{,}5}=\SI{0.893}{\milli\ampere}$.}
\end{questao}

\begin{questao}[3]{}
Um capacitor de $\SI{100}{\micro\farad}$ inicialmente a $\SI{20}{\volt}$ é
descarregado por $R=\SI{1}{\kilo\ohm}$. Mostre, por integração da potência no
resistor, que toda a energia inicialmente armazenada é dissipada e calcule-a.
\resposta{$E=\SI{20}{\milli\joule}$}
\resolucao{$v_C=V_0e^{-t/RC}$ e $i=(V_0/R)e^{-t/RC}$. Logo
$p_R=Ri^2=(V_0^2/R)e^{-2t/RC}$. Integrando de zero a infinito:
\[
E_R=\frac{V_0^2}{R}\int_0^\infty e^{-2t/RC}\,\mathrm{d}t
=\frac{V_0^2}{R}\frac{RC}{2}=\frac12CV_0^2.
\]
Numericamente, $E_R=\tfrac12\pot{100}{-6}\cdot20^2=\SI{20}{\milli\joule}$.}
\end{questao}

\begin{questao}[3]{}
Projete um atraso RC para que uma tensão de $\SI{5}{\volt}$ alcance
$\SI{3.3}{\volt}$ em $\SI{0.5}{\second}$. Use $R=\SI{100}{\kilo\ohm}$ e
considere o capacitor inicialmente descarregado. Determine $C$ e indique um
valor comercial próximo.
\resposta{$C\approx\SI{4.63}{\micro\farad}$; valor comercial: \SI{4.7}{\micro\farad}}
\resolucao{$3{,}3=5(1-e^{-t/RC})$, então
$C=-t/[R\ln(1-3{,}3/5)]$. Substituindo,
$C\approx\SI{4.63}{\micro\farad}$. Um valor comercial próximo é
\SI{4.7}{\micro\farad}.}
\end{questao}

\begin{questao}[3]{}
Deseja-se que a corrente de um circuito RL alcance \SI{95}{\percent} de seu valor
final de $\SI{1}{\ampere}$ em $\SI{20}{\milli\second}$. Se
$L=\SI{50}{\milli\henry}$, determine a resistência e a tensão da fonte.
\resposta{$R\approx\SI{7.49}{\ohm}$; $V\approx\SI{7.49}{\volt}$}
\resolucao{$0{,}95=1-e^{-t/\tau}$ implica
$\tau=-t/\ln0{,}05\approx\SI{6.68}{\milli\second}$. Como $\tau=L/R$,
$R=0{,}05/0{,}00668\approx\SI{7.49}{\ohm}$. Para corrente final de \SI{1}{\ampere},
$V=RI_f\approx\SI{7.49}{\volt}$.}
\end{questao}

\begin{questao}[3]{}
Um RLC em série tem $R=\SI{20}{\ohm}$, $L=\SI{0.1}{\henry}$ e
$C=\SI{100}{\micro\farad}$. Determine os polos naturais.
\resposta{$s_{1,2}=-100\pm j300\ \si{\per\second}$}
\resolucao{$\alpha=R/(2L)=100$ e $\omega_0\approx316{,}2$. Assim
$\omega_d=\sqrt{316{,}2^2-100^2}=300$. Para o caso subamortecido,
$s_{1,2}=-\alpha\pm j\omega_d=-100\pm j300\ \si{\per\second}$.}
\end{questao}

\begin{questao}[3]{}
Uma resposta criticamente amortecida é
$v(t)=(A+Bt)e^{-200t}$. Se $v(0)=\SI{10}{\volt}$ e
$\dot v(0)=-\SI{500}{\volt\per\second}$, determine $A$ e $B$.
\resposta{$A=\SI{10}{\volt}$; $B=\SI{1500}{\volt\per\second}$}
\resolucao{De $v(0)=A$, segue $A=10$. Derivando,
$\dot v(0)=B-200A=-500$. Logo $B=-500+2000=\SI{1500}{\volt\per\second}$.}
\end{questao}

\begin{questao}[3]{}
Para um RLC em série com $R=\SI{100}{\ohm}$, $L=\SI{0.1}{\henry}$ e
$C=\SI{100}{\micro\farad}$, determine os dois polos reais.
\resposta{$s_1\approx-\SI{112.7}{\per\second}$; $s_2\approx-\SI{887.3}{\per\second}$}
\resolucao{A equação característica é
$s^2+(R/L)s+1/(LC)=s^2+1000s+100000=0$. Portanto
\[
s=\frac{-1000\pm\sqrt{1000^2-4\cdot100000}}{2}
=\frac{-1000\pm\sqrt{600000}}{2},
\]
resultando aproximadamente em $-112{,}7$ e $-887{,}3\ \si{\per\second}$.}
\end{questao}

\begin{questao}[3]{}
Determine a resistência crítica de um RLC \textbf{paralelo} com
$L=\SI{50}{\milli\henry}$ e $C=\SI{200}{\micro\farad}$.
\resposta{$R_{\mathrm{crit}}\approx\SI{7.91}{\ohm}$}
\resolucao{No circuito RLC em paralelo, a condição crítica é
$1/(2RC)=1/\sqrt{LC}$. Logo
$R_{\mathrm{crit}}=\tfrac12\sqrt{L/C}
=\tfrac12\sqrt{0{,}05/\pot{200}{-6}}\approx\SI{7.91}{\ohm}$.}
\end{questao}

\begin{questao}[3]{}
Para um sistema subamortecido de segunda ordem com $\zeta=0{,}5$, estime o
sobressinal percentual de pico usando
$M_p=e^{-\pi\zeta/\sqrt{1-\zeta^2}}$.
\resposta{$M_p\approx\SI{16.3}{\percent}$}
\resolucao{Substituindo $\zeta=0{,}5$:
$M_p=e^{-\pi(0{,}5)/\sqrt{0{,}75}}\approx0{,}163$. Portanto o sobressinal é
aproximadamente \SI{16.3}{\percent}.}
\end{questao}

\begin{questao}[3]{}
Um RLC em série possui $R=\SI{20}{\ohm}$, $L=\SI{100}{\milli\henry}$ e
$C=\SI{100}{\micro\farad}$. Calcule o fator de amortecimento $\zeta$.
\resposta{$\zeta\approx0{,}316$}
\resolucao{$\alpha=R/(2L)=100$ e $\omega_0\approx316{,}2$. Assim
$\zeta=\alpha/\omega_0\approx100/316{,}2\approx0{,}316$.}
\end{questao}

\begin{questao}[3]{}
Classifique os quatro casos da tabela.

\begin{table}[H]
\centering\small\sffamily
\begin{tabular}{@{}c c c c c@{}}
\toprule
\textbf{Caso} & \textbf{Topologia} & $R$ & $L$ & $C$\\
\midrule
A & série & \SI{20}{\ohm} & \SI{0.1}{\henry} & \SI{100}{\micro\farad}\\
B & série & \SI{63.25}{\ohm} & \SI{0.1}{\henry} & \SI{100}{\micro\farad}\\
C & série & \SI{100}{\ohm} & \SI{0.1}{\henry} & \SI{100}{\micro\farad}\\
D & paralelo & \SI{20}{\ohm} & \SI{50}{\milli\henry} & \SI{200}{\micro\farad}\\
\bottomrule
\end{tabular}
\end{table}
\resposta{A subamortecido; B aproximadamente crítico; C superamortecido; D subamortecido}
\resolucao{A, B e C possuem $\omega_0\approx316{,}2$. No circuito RLC em série,
$\alpha=R/(0{,}2)$: A dá $100<316{,}2$; B dá aproximadamente $316{,}25$; C dá
$500>316{,}2$. No caso D, $\alpha=1/(2RC)=125$ e
$\omega_0=1/\sqrt{0{,}05\cdot\pot{200}{-6}}\approx316{,}2$, logo também é
subamortecido.}
\end{questao}

\begin{questao}[3]{}
Em certo instante, o capacitor de $C=\SI{100}{\micro\farad}$ apresenta
tensão de $\SI{10}{\volt}$, enquanto o indutor de
$L=\SI{50}{\milli\henry}$ conduz corrente de $\SI{0.2}{\ampere}$. Qual é a energia total armazenada naquele instante?
\resposta{$E=\SI{6}{\milli\joule}$}
\resolucao{$E_C=\tfrac12CV^2=\SI{5}{\milli\joule}$ e
$E_L=\tfrac12LI^2=\tfrac12(0{,}05)(0{,}2)^2=\SI{1}{\milli\joule}$. A soma é
\SI{6}{\milli\joule}.}
\end{questao}

\begin{questao}[3]{}
Antes de uma comutação, um capacitor possui $v_C=\SI{6}{\volt}$ e um indutor
possui $i_L=\SI{0.3}{\ampere}$. Logo após a comutação, ambos passam a integrar
uma nova rede resistiva. Quais condições iniciais devem ser usadas na nova rede?
\resposta{$v_C(0^+)=\SI{6}{\volt}$ e $i_L(0^+)=\SI{0.3}{\ampere}$}
\resolucao{As variáveis de estado associadas aos elementos ideais armazenadores
são contínuas: a tensão do capacitor e a corrente do indutor mantêm seus valores
no instante da comutação. A nova topologia altera suas derivadas e seus valores
finais, não os valores instantâneos em $0^+$.}
\end{questao}

\begin{questao}[3]{}
Para um RLC em série alimentado por um degrau de $\SI{10}{\volt}$, use
$R=\SI{20}{\ohm}$, $L=\SI{0.1}{\henry}$ e $C=\SI{100}{\micro\farad}$. Escreva
a equação diferencial para $v_C(t)$ após a comutação.
\resposta{$\ddot v_C+200\dot v_C+100000v_C=1000000$}
\resolucao{Pela LKT,
$LC\ddot v_C+RC\dot v_C+v_C=V_s$. Como $LC=\pot{1}{-5}$ e
$RC=\pot{2}{-3}$, dividindo por $LC$ resulta
\[
\ddot v_C+200\dot v_C+100000v_C=10\cdot100000=1000000.
\]}
\end{questao}

\begin{questao}[3]{}
No circuito da questão anterior, admita energia inicial nula. Determine
$v_C(0^+)$, $\dot v_C(0^+)$ e $\ddot v_C(0^+)$.
\resposta{$v_C(0^+)=0$; $\dot v_C(0^+)=0$; $\ddot v_C(0^+)=\SI{1e6}{\volt\per\second\squared}$}
\resolucao{Energia inicial nula implica $v_C(0^+)=0$ e $i_L(0^+)=0$. Como
$i_C=C\dot v_C$, segue $\dot v_C(0^+)=0$. Substituindo esses valores na equação
$\ddot v_C+200\dot v_C+100000v_C=1000000$, obtém-se
$\ddot v_C(0^+)=\SI{1e6}{\volt\per\second\squared}$.}
\end{questao}

\gabarito[Nível 3]

% =====================================================================
% CAP. 11 - NIVEL 3 - QUESTOES GRAFICAS E INTEGRADAS
% =====================================================================
\blocoaprofundamento

\begin{questao}[3]{}
Um capacitor de $\SI{40}{\micro\farad}$ permaneceu ligado por muito tempo ao
circuito (a). Em $t=0$, a chave transfere instantaneamente o capacitor para o
circuito (b). Determine $v_C(t)$ para $t>0$, o instante em que a tensão cruza
zero e a corrente no resistor $R_2$ imediatamente após a comutação, tomando como
positivo o sentido do capacitor para a fonte de $-\SI{5}{\volt}$.

\circuitoqq{%
\begin{circuitikz}[scale=0.88,transform shape,line width=0.8pt,every node/.append style={font=\scriptsize}]
\draw (0,0) to[american voltage source,l_={$\SI{15}{\volt}$}] (0,3)
      to[R,l={$R_1=\SI{5}{\kilo\ohm}$}] (3.2,3)
      to[C,v^>={$v_C$}] (3.2,0) -- (0,0);
\node[anchor=east] at (2.95,1.5) {$C=\SI{40}{\micro\farad}$};
\node[font=\sffamily\bfseries\scriptsize] at (1.6,-0.6) {(a) $t<0$};
\begin{scope}[xshift=7.6cm]
\draw (0,0) to[american voltage source,l_={$-\SI{5}{\volt}$}] (0,3)
      to[R,l={$R_2=\SI{10}{\kilo\ohm}$}] (3.2,3)
      to[C,v^>={$v_C$}] (3.2,0) -- (0,0);
\node[anchor=east] at (2.95,1.5) {$C=\SI{40}{\micro\farad}$};
\node[font=\sffamily\bfseries\scriptsize] at (1.6,-0.6) {(b) $t>0$};
\end{scope}
\end{circuitikz}}{Circuitos equivalentes antes e depois da comutação.}

\resposta{$v_C(t)=-5+20e^{-t/0{,}4}\ \si{\volt}$;
$t_{v=0}\approx\SI{0.555}{\second}$; $i_{R_2}(0^+)=\SI{2}{\milli\ampere}$}
\resolucao{Em regime permanente no circuito (a), o capacitor está aberto e não há
queda em $R_1$, de modo que $v_C(0^-)=\SI{15}{\volt}$. Logo
$v_C(0^+)=\SI{15}{\volt}$. Para $t>0$, o valor final é $-\SI{5}{\volt}$ e
\[
\tau=R_2C=(10\,\mathrm{k\Omega})(40\,\mathrm{\mu F})=\SI{0.4}{\second}.
\]
Portanto
\[
v_C(t)=-5+(15+5)e^{-t/0{,}4}=-5+20e^{-t/0{,}4}.
\]
No cruzamento por zero, $e^{-t/0{,}4}=1/4$, logo
$t=0{,}4\ln4\approx\SI{0.555}{\second}$. A corrente inicial em $R_2$ é
$(15-(-5))/10\,\mathrm{k\Omega}=\SI{2}{\milli\ampere}$.}
\end{questao}

\begin{questao}[3]{}
A chave permaneceu fechada por muito tempo e abre em $t=0$. O indutor passa a
liberar sua energia pela associação $R_2\parallel R_3$. Determine a corrente
$i_L(t)$ para $t>0$ e a energia ainda armazenada no indutor em
$t=\SI{50}{\milli\second}$.

\circuitoqq{%
\begin{circuitikz}[scale=0.92,transform shape,line width=0.8pt,every node/.append style={font=\scriptsize}]
\draw (0,0) to[american voltage source,l_={$\SI{30}{\volt}$}] (0,3)
      to[switch] (1.8,3)
      to[R,l={$R_1=\SI{5}{\ohm}$}] (4.0,3) coordinate(a);
\node[anchor=south] at (0.9,3.30) {abre em $t=0$};
\draw (a) to[L,i>^={$i_L$}] (4.0,0) -- (0,0);
\node[anchor=east] at (3.75,1.5) {$L=\SI{200}{\milli\henry}$};
\draw (a) -- (6.3,3) to[R,l_={$R_2=\SI{8}{\ohm}$}] (6.3,0) -- (4.0,0);
\draw (a) -- (8.6,3) to[R,l_={$R_3=\SI{24}{\ohm}$}] (8.6,0) -- (6.3,0);
\end{circuitikz}}{Descarga de um indutor por uma rede resistiva equivalente.}

\resposta{$i_L(t)=6e^{-t/0{,}0333}\ \si{\ampere}$;
$W_L(50\,\mathrm{ms})\approx\SI{0.179}{\joule}$}
\resolucao{Antes da abertura, o indutor ideal está em curto e os resistores
$R_2$ e $R_3$ ficam sem tensão. Assim $i_L(0^-)=30/5=\SI{6}{\ampere}$ e
$i_L(0^+)=\SI{6}{\ampere}$. Para $t>0$,
\[
R_{\mathrm{eq}}=8\parallel24=\SI{6}{\ohm},\qquad
\tau=\frac{0{,}2}{6}\approx\SI{33.3}{\milli\second}.
\]
Logo $i_L(t)=6e^{-t/\tau}$. Em $\SI{50}{\milli\second}$,
$i_L=6e^{-1{,}5}\approx\SI{1.339}{\ampere}$. A energia remanescente é
\[
W_L=\frac12Li_L^2\approx\frac12(0{,}2)(1{,}339)^2
\approx\SI{0.179}{\joule}.
\]}
\end{questao}

\begin{questao}[3]{}
O capacitor possui $v_C(0^-)=\SI{10}{\volt}$. Em $t=0$ a chave fecha e o conecta
à rede mostrada. Determine $v_C(t)$, a resistência equivalente vista pelo
capacitor e o tempo necessário para que o erro em relação ao valor final seja
menor que \SI{1}{\percent} da variação inicial.

\circuitoqq{%
\begin{circuitikz}[scale=0.91,transform shape,line width=0.8pt,every node/.append style={font=\scriptsize}]
\draw (0,0) to[american voltage source,l_={$\SI{18}{\volt}$}] (0,3)
      to[R,l={$R_1=\SI{6}{\kilo\ohm}$}] (3.2,3) coordinate(b);
\draw (b) to[R,l_={$R_2=\SI{3}{\kilo\ohm}$}] (3.2,0) -- (0,0);
\draw (b) to[R,l={$R_3=\SI{2}{\kilo\ohm}$}] (5.7,3)
      to[switch] (7.2,3)
      to[C,v^>={$v_C$}] (7.2,0) -- (3.2,0);
\node[anchor=south] at (6.45,3.30) {fecha em $t=0$};
\node[anchor=east] at (6.95,1.5) {$C=\SI{25}{\micro\farad}$};
\end{circuitikz}}{Capacitor conectado a uma rede de Thévenin não trivial.}

\resposta{$R_{\mathrm{eq}}=\SI{4}{\kilo\ohm}$; $\tau=\SI{0.10}{\second}$;
$v_C(t)=6+4e^{-10t}\ \si{\volt}$; $t_{1\%}\approx\SI{0.461}{\second}$}
\resolucao{Em regime final não circula corrente em $R_3$, portanto o capacitor
assume a tensão do divisor formado por $R_1$ e $R_2$:
$V_f=18\cdot3/(6+3)=\SI{6}{\volt}$. Para obter a resistência vista pelo capacitor,
a fonte ideal de tensão é curto-circuitada:
\[
R_{\mathrm{eq}}=R_3+(R_1\parallel R_2)=2\,\mathrm{k\Omega}+2\,\mathrm{k\Omega}
=\SI{4}{\kilo\ohm}.
\]
Logo $\tau=R_{\mathrm{eq}}C=\SI{0.10}{\second}$ e
$v_C(t)=6+(10-6)e^{-t/0{,}1}=6+4e^{-10t}$. O erro normalizado é
$e^{-t/\tau}$. Para ficar abaixo de \SI{1}{\percent},
$t>\tau\ln100\approx4{,}605\tau=\SI{0.461}{\second}$.}
\end{questao}

\begin{questao}[3]{}
Para $t<0$, a fonte de corrente de $\SI{2}{\ampere}$ permaneceu conectada por
muito tempo ao RLC em paralelo. Em $t=0$ a chave abre e a fonte é removida. Determine
o tipo de resposta natural e obtenha $v(t)$ para $t>0$, considerando a polaridade
indicada.

\circuitoqq{%
\begin{circuitikz}[scale=0.94,transform shape,line width=0.8pt,every node/.append style={font=\scriptsize}]
\draw (0,0) to[american current source] (0,3)
      to[switch] (1.9,3) -- (7.3,3);
\node[anchor=east] at (-0.25,1.5) {$\SI{2}{\ampere}$};
\node[anchor=south] at (0.95,3.30) {abre em $t=0$};
\draw (2.6,3) to[R] (2.6,0);
\node[anchor=east] at (2.35,1.5) {$R=\SI{10}{\ohm}$};
\draw (4.9,3) to[C,v^>={$v$}] (4.9,0);
\node[anchor=east] at (4.65,1.5) {$C=\SI{200}{\micro\farad}$};
\draw (7.3,3) to[L,i>^={$i_L$}] (7.3,0) -- (0,0);
\node[anchor=east] at (7.05,1.5) {$L=\SI{50}{\milli\henry}$};
\end{circuitikz}}{Resposta natural de um circuito RLC em paralelo com energia inicialmente no indutor.}

\resposta{Subamortecida;
$v(t)\approx-51{,}64e^{-250t}\sin(193{,}65t)\ \si{\volt}$}
\resolucao{Antes da abertura, em regime CC, o capacitor é aberto e o indutor é
curto. Assim $v(0^-)=0$ e $i_L(0^-)=\SI{2}{\ampere}$; por continuidade,
$v(0^+)=0$ e $i_L(0^+)=\SI{2}{\ampere}$. Para o RLC em paralelo,
\[
\alpha=\frac{1}{2RC}=\SI{250}{\per\second},\qquad
\omega_0=\frac{1}{\sqrt{LC}}\approx\SI{316.23}{\radian\per\second},
\]
logo a resposta é subamortecida e
$\omega_d\approx\SI{193.65}{\radian\per\second}$. Pela KCL em $0^+$,
$C\dot v(0^+)+i_L(0^+)=0$, portanto
$\dot v(0^+)=-2/(200\,\mathrm{\mu F})=-\SI{10000}{\volt\per\second}$.
Como $v(0)=0$, a parcela cossenoidal é nula e
$B=\dot v(0)/\omega_d\approx-51{,}64$, resultando na expressão indicada.}
\end{questao}

\begin{questao}[3]{}
O circuito RLC em série parte sem energia armazenada e a chave fecha em $t=0$.
Determine a resposta completa $v_C(t)$ e o valor do primeiro pico da tensão do
capacitor.

\circuitoqq{%
\begin{circuitikz}[scale=0.90,transform shape,line width=0.8pt,every node/.append style={font=\scriptsize}]
\draw (0,0) to[american voltage source,l_={$\SI{10}{\volt}$}] (0,3)
      to[switch] (1.8,3)
      to[R,l={$R=\SI{20}{\ohm}$}] (4.0,3)
      to[L,l={$L=\SI{100}{\milli\henry}$}] (6.5,3)
      to[C,v^>={$v_C$}] (9.0,3)
      -- (9.0,0) -- (0,0);
\node[anchor=south] at (0.9,3.30) {fecha em $t=0$};
\node[anchor=south] at (7.75,3.30) {$C=\SI{100}{\micro\farad}$};
\end{circuitikz}}{Resposta ao degrau de um circuito RLC em série subamortecido.}

\resposta{$v_C(t)=10\left[1-e^{-100t}\left(\cos300t+\frac13\sin300t\right)\right]\,\si{\volt}$;
$v_{C,\mathrm{pico}}\approx\SI{13.51}{\volt}$ em $t_p\approx\SI{10.47}{\milli\second}$}
\resolucao{Temos
$\alpha=R/(2L)=\SI{100}{\per\second}$,
$\omega_0\approx\SI{316.23}{\radian\per\second}$ e
$\omega_d=\SI{300}{\radian\per\second}$. Para entrada degrau e condições iniciais
nulas, a resposta do capacitor é
\[
v_C(t)=V_s\left[1-e^{-\alpha t}\left(\cos\omega_dt+
\frac{\alpha}{\omega_d}\sin\omega_dt\right)\right].
\]
Substituindo os valores, obtém-se a expressão indicada. O primeiro pico ocorre em
$t_p=\pi/\omega_d\approx\SI{10.47}{\milli\second}$. Nesse instante,
$\sin(\omega_dt_p)=0$ e $\cos(\omega_dt_p)=-1$, logo
$v_C(t_p)=10[1+e^{-100t_p}]\approx\SI{13.51}{\volt}$.}
\end{questao}

\begin{questao}[3]{}
Em laboratório, mediu-se a corrente natural do RLC em série mostrado. A distância
entre picos consecutivos é aproximadamente $T_d=\SI{8}{\milli\second}$ e a
envoltória cai à metade a cada $\SI{16}{\milli\second}$. Sabendo que
$R=\SI{8}{\ohm}$, estime $\alpha$, $\omega_d$, $L$ e $C$.

\circuitoqq{%
\begin{circuitikz}[scale=0.96,transform shape,line width=0.8pt,every node/.append style={font=\scriptsize}]
\draw (0,0) to[R] (2.2,0)
      to[L] (4.5,0)
      to[C] (6.8,0) -- (6.8,-2) -- (0,-2) -- (0,0);
\node[anchor=south] at (1.1,0.25) {$R=\SI{8}{\ohm}$};
\node[anchor=south] at (3.35,0.25) {$L$};
\node[anchor=south] at (5.65,0.25) {$C$};
\draw[->] (0.9,0.70) -- (2.2,0.70);
\node[anchor=south] at (1.55,0.82) {$i(t)$};
\end{circuitikz}}{RLC em série identificado a partir de uma oscilação medida.}

\begin{figure}[H]
\centering
\begin{tikzpicture}
\begin{axis}[
 width=.88\linewidth,height=5.2cm,
 xlabel={$t$ (ms)},ylabel={$i(t)$ (u.a.)},xmin=0,xmax=32,ymin=-2.2,ymax=2.2,
 grid=both,samples=300,
 tick label style={font=\scriptsize},label style={font=\small\sffamily}]
\addplot[thick,azulprof,domain=0:32] {2*exp(-0.0433217*x)*sin(deg(0.785398*x))};
\addplot[dashed,laranja,domain=0:32] {2*exp(-0.0433217*x)};
\addplot[dashed,laranja,domain=0:32] {-2*exp(-0.0433217*x)};
\end{axis}
\end{tikzpicture}
\caption{Oscilação amortecida medida e envoltória exponencial.}
\end{figure}

\resposta{$\alpha\approx\SI{43.3}{\per\second}$;
$\omega_d\approx\SI{785.4}{\radian\per\second}$;
$L\approx\SI{92.3}{\milli\henry}$; $C\approx\SI{17.5}{\micro\farad}$}
\resolucao{A envoltória obedece $e^{-\alpha t}$. Como ela cai à metade em
$\SI{16}{\milli\second}$,
\[
\alpha=\frac{\ln2}{0{,}016}\approx\SI{43.3}{\per\second}.
\]
O período amortecido fornece
$\omega_d=2\pi/T_d\approx\SI{785.4}{\radian\per\second}$. Para RLC em série,
$\alpha=R/(2L)$, então
$L=R/(2\alpha)\approx\SI{92.3}{\milli\henry}$. Além disso,
$\omega_0=\sqrt{\omega_d^2+\alpha^2}\approx\SI{786.6}{\radian\per\second}$ e
\[
C=\frac{1}{L\omega_0^2}\approx\SI{17.5}{\micro\farad}.
\]}
\end{questao}

\gabarito[Nível 3 --- análise integrada e formas de onda]

% =====================================================================
% CAP. 11 - NIVEL 4
% =====================================================================
\blocodesafio

\begin{questao}[4]{}
Em laboratório, um circuito RC com $R=\SI{5}{\kilo\ohm}$ e
$C=\SI{10}{\micro\farad}$ deveria apresentar $\tau=\SI{50}{\milli\second}$.
Entretanto, o ajuste exponencial do osciloscópio fornece
$\tau_{\mathrm{med}}=\SI{80}{\milli\second}$. A capacitância medida está
próxima do valor nominal. Estime a resistência efetiva vista pelo capacitor e proponha
uma causa elétrica plausível para a discrepância.
\resposta{$R_{\mathrm{eq}}\approx\SI{8}{\kilo\ohm}$; há cerca de \SI{3}{\kilo\ohm}
adicionais na trajetória equivalente}
\resolucao{$R_{\mathrm{eq}}=\tau/C=0{,}08/(\pot{10}{-6})=\SI{8}{\kilo\ohm}$.
Como o resistor nominal é de \SI{5}{\kilo\ohm}, há uma resistência
adicional de aproximadamente \SI{3}{\kilo\ohm}. Essa diferença pode decorrer da resistência de saída da fonte, de uma
rede de polarização não considerada, de um resistor em série, de uma chave ou
sensor, ou de uma topologia
que não foi corretamente reduzida ao equivalente de Thévenin.}
\end{questao}

\begin{questao}[4]{}
Um capacitor de $\SI{470}{\micro\farad}$ será carregado por uma fonte de
$\SI{24}{\volt}$. Deseja-se limitar a corrente inicial a
$\SI{100}{\milli\ampere}$ e atingir pelo menos \SI{95}{\percent} da tensão final
em menos de $\SI{0.5}{\second}$. Projete o resistor mínimo em série, verifique o
tempo de carga e estime a energia total dissipada nele durante uma carga completa.
\resposta{$R=\SI{240}{\ohm}$; $t_{95}\approx\SI{0.338}{\second}$; $E_R\approx\SI{0.135}{\joule}$}
\resolucao{Pela corrente inicial, $R\ge24/0{,}1=\SI{240}{\ohm}$. Com esse valor,
$\tau=RC=240\cdot\pot{470}{-6}=\SI{112.8}{\milli\second}$. Para \SI{95}{\percent},
$t=-\tau\ln0{,}05\approx\SI{0.338}{\second}$, atendendo à especificação. A
energia dissipada no resistor durante a carga ideal por fonte de tensão é igual à
energia final do capacitor:
$E_R=\tfrac12CV^2\approx\SI{0.135}{\joule}$. A potência inicial é
$V^2/R=\SI{2.4}{\watt}$, importante para a escolha da tecnologia do resistor.}
\end{questao}

\begin{questao}[4]{}
Uma bobina de relé tem $L=\SI{120}{\milli\henry}$ e resistência de cobre
$\SI{60}{\ohm}$, alimentada em $\SI{24}{\volt}$ até o regime permanente.
Compare duas estratégias de desligamento: (A) diodo de roda livre ideal; (B)
resistor de clamp de $\SI{180}{\ohm}$ no caminho de descarga. Determine a energia
inicial, as constantes de tempo aproximadas e a tensão inicial no resistor de
clamp. Explique o compromisso de projeto.
\resposta{$I_0=\SI{0.4}{\ampere}$; $E_L=\SI{9.6}{\milli\joule}$; $\tau_A\approx\SI{2}{\milli\second}$;
$\tau_B\approx\SI{0.5}{\milli\second}$; $v_{R,\mathrm{clamp}}(0^+)\approx\SI{72}{\volt}$}
\resolucao{Em regime, $I_0=24/60=\SI{0.4}{\ampere}$ e
$E_L=\tfrac12LI_0^2=\SI{9.6}{\milli\joule}$. Com o diodo ideal, a resistência no caminho de descarga é aproximadamente
a resistência do enrolamento da bobina: $\tau_A=L/60=\SI{2}{\milli\second}$. Com o
resistor adicional, $R_{\mathrm{tot}}=60+180=\SI{240}{\ohm}$ e
$\tau_B=L/240=\SI{0.5}{\milli\second}$. A tensão inicial no resistor externo é
$0{,}4\cdot180=\SI{72}{\volt}$. O clamp resistivo desliga o relé mais rápido,
mas produz tensão maior; o diodo reduz a sobretensão, porém prolonga a corrente.}
\end{questao}

\begin{questao}[4]{}
Uma oscilação amortecida medida em um circuito RLC apresenta período entre picos
de aproximadamente $\SI{4}{\milli\second}$ e a envoltória cai à metade a cada
$\SI{10}{\milli\second}$. O capacitor é $C=\SI{10}{\micro\farad}$. Estime
$\alpha$, $\omega_d$, $\omega_0$, $\zeta$, $L$ e a resistência equivalente em série.

\begin{figure}[H]
\centering
\begin{tikzpicture}
\begin{axis}[
 width=0.88\linewidth,height=5.2cm,
 xlabel={$t$ (ms)},ylabel={$v(t)$},xmin=0,xmax=30,
 axis lines=middle,grid=both,samples=300,
 tick label style={font=\scriptsize},label style={font=\small\sffamily}]
\addplot[thick,azulprof,domain=0:30] {10*exp(-0.0693147*x)*cos(deg(1.5707963*x))};
\addplot[dashed,laranja,domain=0:30] {10*exp(-0.0693147*x)};
\addplot[dashed,laranja,domain=0:30] {-10*exp(-0.0693147*x)};
\end{axis}
\end{tikzpicture}
\caption{Exemplo de oscilação amortecida e sua envoltória exponencial.}
\end{figure}

\resposta{$\alpha\approx\SI{69.3}{\per\second}$; $\omega_d\approx\SI{1571}{\radian\per\second}$;
$\omega_0\approx\SI{1572}{\radian\per\second}$; $\zeta\approx0{,}044$;
$L\approx\SI{40.45}{\milli\henry}$; $R\approx\SI{5.61}{\ohm}$}
\resolucao{Da envoltória, $e^{-\alpha(0{,}01)}=1/2$, então
$\alpha=\ln2/0{,}01\approx\SI{69.3}{\per\second}$. Do período,
$\omega_d=2\pi/0{,}004\approx\SI{1571}{\radian\per\second}$. Logo
$\omega_0=\sqrt{\omega_d^2+\alpha^2}\approx\SI{1572}{\radian\per\second}$ e
$\zeta\approx0{,}044$. Como $\omega_0^2=1/(LC)$,
$L\approx1/(C\omega_0^2)\approx\SI{40.45}{\milli\henry}$. Por fim,
$R=2\alpha L\approx\SI{5.61}{\ohm}$.}
\end{questao}

\begin{questao}[4]{}
Um estágio RLC em série usa $L=\SI{20}{\milli\henry}$ e
$C=\SI{10}{\micro\farad}$. O requisito é obter a resposta mais rápida possível
sem oscilação. Determine o resistor ideal e compare qualitativamente o uso de
$\SI{45}{\ohm}$ e $\SI{180}{\ohm}$.
\resposta{$R_{\mathrm{crit}}\approx\SI{89.4}{\ohm}$; \SI{45}{\ohm} é subamortecido;
\SI{180}{\ohm} é superamortecido}
\resolucao{$R_{\mathrm{crit}}=2\sqrt{L/C}=2\sqrt{0{,}02/\pot{10}{-6}}
\approx\SI{89.4}{\ohm}$. Com \SI{45}{\ohm}, $R<R_{\mathrm{crit}}$ e haverá
oscilações amortecidas (ringing) e sobressinal. Com \SI{180}{\ohm}, $R>R_{\mathrm{crit}}$ e a resposta será
monotônica, porém mais lenta. O amortecimento crítico é o limite de resposta
monotônica mais rápida para o modelo ideal.}
\end{questao}

\begin{questao}[4]{}
A tabela mostra a resposta de um RC inicialmente descarregado cuja tensão final é
$\SI{10}{\volt}$.

\begin{table}[H]
\centering\small\sffamily
\begin{tabular}{@{}c c@{}}
\toprule
$t$ & $v_C(t)$\\
\midrule
\SI{20}{\milli\second} & \SI{6.32}{\volt}\\
\SI{40}{\milli\second} & \SI{8.65}{\volt}\\
\bottomrule
\end{tabular}
\end{table}

Verifique se os dados são compatíveis com uma única constante de tempo e estime
$R$ se $C=\SI{4.7}{\micro\farad}$.
\resposta{$\tau\approx\SI{20}{\milli\second}$; $R\approx\SI{4.26}{\kilo\ohm}$}
\resolucao{Em $t=\tau$, uma carga RC atinge \SI{63.2}{\percent}, exatamente o
primeiro ponto. Em $2\tau$, atinge $1-e^{-2}=0{,}8647$, coerente com
\SI{8.65}{\volt}. Logo $\tau\approx\SI{20}{\milli\second}$. Assim
$R=\tau/C=0{,}02/(\pot{4.7}{-6})\approx\SI{4.26}{\kilo\ohm}$.}
\end{questao}

\begin{questao}[4]{}
Após uma comutação, um circuito RLC em série fica sem fonte independente e possui
$R=\SI{20}{\ohm}$, $L=\SI{50}{\milli\henry}$ e
$C=\SI{100}{\micro\farad}$. Adote $v_C(0^+)=\SI{5}{\volt}$ e
$i(0^+)=\SI{0.1}{\ampere}$ entrando no terminal positivo do capacitor. Determine
a expressão de $v_C(t)$.
\resposta{$v_C(t)=e^{-200t}[5\cos(400t)+5\sin(400t)]\ \si{\volt}$}
\resolucao{$\alpha=R/(2L)=200$, $\omega_0=1/\sqrt{LC}\approx447{,}2$ e
$\omega_d=\sqrt{\omega_0^2-\alpha^2}=400$. Portanto
$v_C=e^{-200t}(A\cos400t+B\sin400t)$. De $v_C(0)=5$, $A=5$. Como
$i=C\dot v_C$, $\dot v_C(0)=0{,}1/(\pot{100}{-6})=\SI{1000}{\volt\per\second}$.
Então $\dot v_C(0)=-200A+400B=1000$, resultando em $B=5$.}
\end{questao}

\begin{questao}[4]{}
Projete uma rede de temporização para uma entrada de $\SI{3.3}{\volt}$ e
$C=\SI{1}{\micro\farad}$. Na carga, a tensão deve alcançar
$\SI{2}{\volt}$ em aproximadamente $\SI{8}{\milli\second}$. No desligamento,
deve cair de $\SI{3.3}{\volt}$ para menos de $\SI{0.5}{\volt}$ em
$\SI{5}{\milli\second}$. Mostre que um único resistor não atende bem às duas
metas e proponha resistores independentes de carga e descarga usando um diodo.
\resposta{$R_{\mathrm{carga}}\approx\SI{8.59}{\kilo\ohm}$; $R_{\mathrm{descarga}}\lesssim\SI{2.65}{\kilo\ohm}$;
valores práticos: \SI{8.2}{\kilo\ohm} e \SI{2.4}{\kilo\ohm}}
\resolucao{Na carga,
$2=3{,}3(1-e^{-8\mathrm{ms}/RC})$, fornecendo
$R\approx\SI{8.59}{\kilo\ohm}$. Se esse mesmo resistor fosse usado na descarga,
o tempo até \SI{0.5}{\volt} seria $RC\ln(3{,}3/0{,}5)\approx\SI{16.2}{\milli\second}$,
acima do limite. Para a descarga em no máximo \SI{5}{\milli\second},
$R\le t/[C\ln(3{,}3/0{,}5)]\approx\SI{2.65}{\kilo\ohm}$. Um diodo permite
separar os caminhos; \SI{8.2}{\kilo\ohm} e \SI{2.4}{\kilo\ohm} são escolhas
comerciais coerentes.}
\end{questao}

\begin{questao}[4]{}
Para $L=\SI{20}{\milli\henry}$ e $C=\SI{10}{\micro\farad}$, compare os três
resistores da tabela e escolha o mais adequado quando se deseja resposta rápida
sem sobressinal.

\begin{table}[H]
\centering\small\sffamily
\begin{tabular}{@{}c c c@{}}
\toprule
\textbf{Opção} & $R$ & \textbf{Observação a determinar}\\
\midrule
A & \SI{40}{\ohm} & ?\\
B & \SI{89.4}{\ohm} & ?\\
C & \SI{180}{\ohm} & ?\\
\bottomrule
\end{tabular}
\end{table}

\resposta{A: $\zeta\approx0{,}447$, subamortecido; B: crítico; C: $\zeta\approx2{,}01$, superamortecido; escolher B}
\resolucao{Como $R_{\mathrm{crit}}\approx\SI{89.4}{\ohm}$, no RLC em série vale
$\zeta=R/R_{\mathrm{crit}}$. Assim A fornece $\zeta\approx0{,}447$ e oscila; B
fica praticamente em $\zeta=1$; C fornece $\zeta\approx2{,}01$ e é lento por
superamortecimento. Para rapidez sem sobressinal, B é a escolha adequada.}
\end{questao}

\begin{questao}[4]{}
Um projeto prevê $R=\SI{20}{\ohm}$, $L=\SI{0.1}{\henry}$ e
$C=\SI{100}{\micro\farad}$, portanto deveria ser subamortecido. Contudo, o ajuste
da resposta medida indica dois polos reais em aproximadamente
$-\SI{112.7}{\per\second}$ e $-\SI{887.3}{\per\second}$. Admitindo $L$ e $C$
corretos, estime a resistência efetiva em série e diagnostique a discrepância.
\resposta{$R_{\mathrm{efetivo}}\approx\SI{100}{\ohm}$; há cerca de \SI{80}{\ohm}
de resistência em série não contabilizada}
\resolucao{Para $s^2+(R/L)s+1/(LC)=0$, a soma dos polos é $-R/L$. A soma medida é
$-1000\ \si{\per\second}$; com $L=\SI{0.1}{\henry}$,
$R=1000L=\SI{100}{\ohm}$. O produto dos polos é aproximadamente
$100000\ \si{\per\second\squared}$, coerente com $1/(LC)$, confirmando que $L$ e
$C$ são plausíveis. Portanto existe cerca de \SI{80}{\ohm} extra, por ESR,
resistência do indutor, resistor de proteção, fonte ou fiação.}
\end{questao}

\begin{questao}[4]{}
Um barramento CC de $\SI{400}{\volt}$ possui capacitor de
$\SI{470}{\micro\farad}$. Por segurança, após desligar a fonte deseja-se reduzir
a tensão para no máximo $\SI{50}{\volt}$ em $\SI{60}{\second}$. Dimensione o
maior resistor de descarga admissível, escolha um valor comercial abaixo dele,
estime a potência inicial e a energia que será dissipada.
\resposta{$R_{\max}\approx\SI{61.4}{\kilo\ohm}$; escolha \SI{56}{\kilo\ohm};
$t\approx\SI{54.7}{\second}$; $P_0\approx\SI{2.86}{\watt}$; $E\approx\SI{37.6}{\joule}$}
\resolucao{Na descarga, $50=400e^{-t/RC}$. Assim
$R=t/[C\ln(400/50)]\approx\SI{61.4}{\kilo\ohm}$. Um valor de
\SI{56}{\kilo\ohm} fornece $t=RC\ln8\approx\SI{54.7}{\second}$. A potência
inicial é $P_0=V^2/R=400^2/56000\approx\SI{2.86}{\watt}$, portanto é prudente
usar resistor com margem de potência e tensão. A energia inicialmente armazenada,
e posteriormente dissipada, é $E=\tfrac12CV^2\approx\SI{37.6}{\joule}$.}
\end{questao}

\gabarito[Nível 4]

% =====================================================================
% CAP. 11 - NIVEL 4 - PROJETO, DIAGNOSTICO E CIRCUITOS INTEGRADOS
% =====================================================================
\blocodesafio

\begin{questao}[4]{}
Um sensor pode ser modelado por uma fonte de Thévenin de $\SI{5}{\volt}$ com
resistência interna $R_s=\SI{1}{\kilo\ohm}$. Ao fechar a chave em $t=0$, ele
alimenta uma entrada com $R_L=\SI{9}{\kilo\ohm}$ em paralelo com um capacitor
$C$, inicialmente descarregado. O sistema digital reconhece nível alto quando
$v_C\ge\SI{3}{\volt}$ e exige que isso ocorra em no máximo
$\SI{20}{\milli\second}$.

\circuitoqq{%
\begin{circuitikz}[scale=0.90,transform shape,line width=0.8pt,every node/.append style={font=\scriptsize}]
\draw (0,0) to[american voltage source,l_={$\SI{5}{\volt}$}] (0,3)
      to[R,l={$R_s=\SI{1}{\kilo\ohm}$}] (2.5,3)
      to[switch] (4.3,3) -- (6.8,3) coordinate(a);
\node[anchor=south] at (3.4,3.30) {fecha em $t=0$};
\draw (a) to[R,l_={$R_L=\SI{9}{\kilo\ohm}$}] (6.8,0) -- (0,0);
\draw (a) -- (9.2,3) to[C,v^>={$v_C$}] (9.2,0) -- (6.8,0);
\node[anchor=east] at (8.95,1.5) {$C$};
\end{circuitikz}}{Entrada RC carregada por uma fonte com resistência interna e carga finita.}

\begin{itensq}
\item Determine o valor final de $v_C$ e a resistência de Thévenin vista pelo capacitor.
\item Calcule o maior valor de $C$ que atende ao tempo especificado.
\item Escolha entre $\SI{18}{\micro\farad}$ e $\SI{22}{\micro\farad}$ e verifique o tempo de cruzamento de \SI{3}{\volt}.
\item Determine a corrente fornecida pela fonte em $0^+$ e em regime permanente.
\end{itensq}

\resposta{$V_f=\SI{4.5}{\volt}$; $R_{\mathrm{Th}}=\SI{900}{\ohm}$;
$C_{\max}\approx\SI{20.2}{\micro\farad}$; escolha: \SI{18}{\micro\farad};
$t_{3V}\approx\SI{17.8}{\milli\second}$;
$i_s(0^+)=\SI{5}{\milli\ampere}$; $i_s(\infty)=\SI{0.5}{\milli\ampere}$}
\resolucao{Em regime permanente o capacitor é aberto, então
\[
V_f=5\frac{9}{1+9}=\SI{4.5}{\volt}.
\]
Suprimindo a fonte, a resistência vista pelo capacitor é
$R_{\mathrm{Th}}=R_s\parallel R_L=1\parallel9=\SI{0.9}{\kilo\ohm}$.
Como o capacitor parte descarregado,
$v_C(t)=4{,}5(1-e^{-t/(R_{\mathrm{Th}}C)})$. A condição limite é
$3=4{,}5(1-e^{-0{,}02/(900C)})$, ou seja $e^{-0{,}02/(900C)}=1/3$. Logo
\[
C_{\max}=\frac{0{,}02}{900\ln3}\approx\SI{20.2}{\micro\farad}.
\]
O valor de \SI{22}{\micro\farad} viola a especificação; com
$C=\SI{18}{\micro\farad}$, $\tau=\SI{16.2}{\milli\second}$ e
$t_{3V}=\tau\ln3\approx\SI{17.8}{\milli\second}$. Em $0^+$ o capacitor se
comporta como curto, logo $i_s=5/1\,\mathrm{k\Omega}=\SI{5}{\milli\ampere}$.
Em regime, a corrente é $5/(1+9)\,\mathrm{k\Omega}=\SI{0.5}{\milli\ampere}$.}
\end{questao}

\begin{questao}[4]{}
Um barramento CC possui um capacitor de $\SI{470}{\micro\farad}$ carregado a
$\SI{325}{\volt}$. Quando a alimentação é desconectada, um resistor de descarga
$R_b$ deve reduzir a tensão para no máximo $\SI{50}{\volt}$ em
$\SI{60}{\second}$. Despreze outras cargas.

\circuitoqq{%
\begin{circuitikz}[scale=0.94,transform shape,line width=0.8pt,every node/.append style={font=\scriptsize}]
\draw (0,0) to[american voltage source,l_={$\SI{325}{\volt}$}] (0,3)
      to[switch] (2.0,3) -- (5.2,3) coordinate(a);
\node[anchor=south] at (1.0,3.30) {desliga em $t=0$};
\draw (a) to[C,v^>={$v_C$}] (5.2,0) -- (0,0);
\node[anchor=east] at (4.95,1.5) {$C=\SI{470}{\micro\farad}$};
\draw (a) -- (7.8,3) to[R,l_={$R_b$}] (7.8,0) -- (5.2,0);
\end{circuitikz}}{Resistor de descarga de um barramento capacitivo.}

\begin{itensq}
\item Determine o maior valor permitido de $R_b$.
\item Escolha um valor comercial de $\SI{68}{\kilo\ohm}$ e verifique o tempo até \SI{50}{\volt}.
\item Calcule a potência instantânea inicial no resistor e a energia total que ele dissipará.
\end{itensq}

\resposta{$R_{b,\max}\approx\SI{68.2}{\kilo\ohm}$; com \SI{68}{\kilo\ohm},
$t_{50V}\approx\SI{59.8}{\second}$; $p_R(0^+)\approx\SI{1.55}{\watt}$;
$E\approx\SI{24.8}{\joule}$}
\resolucao{Na descarga, $v_C(t)=V_0e^{-t/(R_bC)}$. Impor
$50=325e^{-60/(R_bC)}$ fornece
\[
R_b=\frac{-60}{C\ln(50/325)}\approx\SI{68.2}{\kilo\ohm}.
\]
Assim, \SI{68}{\kilo\ohm} atende à especificação. Com esse valor,
$t=R_bC\ln(325/50)\approx\SI{59.8}{\second}$. A potência inicial é
$p_R(0^+)=V_0^2/R_b\approx\SI{1.55}{\watt}$. Toda a energia inicialmente
armazenada no capacitor é dissipada no resistor:
\[
E=\frac12CV_0^2
=\frac12(470\times10^{-6})(325)^2\approx\SI{24.8}{\joule}.
\]
O projeto real deve ainda considerar a tensão máxima do resistor e sua capacidade
de suportar a energia do pulso de descarga.}
\end{questao}

\begin{questao}[4]{}
Projete o RLC em série abaixo para que a resposta do capacitor a um degrau seja
subamortecida com $\zeta=0{,}5$ e tempo de acomodação de \SI{2}{\percent} não
superior a $\SI{40}{\milli\second}$. Adote
$C=\SI{100}{\micro\farad}$ e use a aproximação
$t_s\approx4/(\zeta\omega_0)$. Escolha o menor $\omega_0$ que atende à
especificação. Determine $L$, $R$, o sobressinal percentual e o instante do
primeiro pico.

\circuitoqq{%
\begin{circuitikz}[scale=0.90,transform shape,line width=0.8pt,every node/.append style={font=\scriptsize}]
\draw (0,0) to[american voltage source,l_={$V_su(t)$}] (0,3)
      to[switch] (1.8,3)
      to[R,l={$R$}] (4.0,3)
      to[L,l={$L$}] (6.5,3)
      to[C,v^>={$v_C$}] (9.0,3)
      -- (9.0,0) -- (0,0);
\node[anchor=south] at (0.9,3.30) {fecha em $t=0$};
\node[anchor=south] at (7.75,3.30) {$C=\SI{100}{\micro\farad}$};
\end{circuitikz}}{Projeto de um circuito RLC em série a partir de especificações transitórias.}

\resposta{$\omega_0=\SI{200}{\radian\per\second}$;
$L=\SI{0.25}{\henry}$; $R=\SI{50}{\ohm}$;
$M_p\approx\SI{16.3}{\percent}$; $t_p\approx\SI{18.1}{\milli\second}$}
\resolucao{Usando a especificação no limite,
\[
0{,}040=\frac{4}{0{,}5\omega_0}
\quad\Rightarrow\quad
\omega_0=\SI{200}{\radian\per\second}.
\]
Como $\omega_0=1/\sqrt{LC}$,
$L=1/(C\omega_0^2)=1/(100\times10^{-6}\cdot200^2)=\SI{0.25}{\henry}$.
Para RLC em série, $\zeta=R/(2L\omega_0)$; portanto
$R=2\zeta L\omega_0=\SI{50}{\ohm}$. O sobressinal é
\[
M_p=e^{-\pi\zeta/\sqrt{1-\zeta^2}}\approx0{,}163,
\]
ou \SI{16.3}{\percent}. A frequência amortecida é
$\omega_d=\omega_0\sqrt{1-\zeta^2}\approx\SI{173.2}{\radian\per\second}$ e
$t_p=\pi/\omega_d\approx\SI{18.1}{\milli\second}$.}
\end{questao}

\begin{questao}[4]{}
O circuito RC em escada abaixo possui \textbf{dois capacitores independentes} e,
portanto, é de segunda ordem. Para $v_s(t)=10u(t)\,\si{\volt}$ e condições
iniciais nulas, obtenha as equações de estado em $v_1$ e $v_2$, determine os polos
naturais e encontre $v_2(t)$.

\circuitoqq{%
\begin{circuitikz}[scale=0.94,transform shape,line width=0.8pt,every node/.append style={font=\scriptsize}]
\draw (0,0) to[american voltage source] (0,3)
      to[R,l={$R_1=\SI{1}{\kilo\ohm}$}] (3.2,3) coordinate(a)
      to[R,l={$R_2=\SI{2}{\kilo\ohm}$}] (6.5,3) coordinate(b);
\node[anchor=east] at (-0.25,1.5) {$10u(t)\,\si{\volt}$};
\draw (a) to[C,v^>={$v_1$}] (3.2,0);
\node[anchor=east] at (2.95,1.5) {$C_1=\SI{100}{\micro\farad}$};
\draw (b) to[C,v^>={$v_2$}] (6.5,0) -- (0,0);
\node[anchor=east] at (6.25,1.5) {$C_2=\SI{50}{\micro\farad}$};
\end{circuitikz}}{Rede RC de segunda ordem com dois elementos armazenadores do mesmo tipo.}

\resposta{$\dot v_1=100u(t)-15v_1+5v_2$;
$\dot v_2=10v_1-10v_2$; polos $-5$ e $-20\ \si{\per\second}$;
$v_2(t)=10-\frac{40}{3}e^{-5t}+\frac{10}{3}e^{-20t}\ \si{\volt}$}
\resolucao{Aplicando KCL aos dois nós,
\[
C_1\dot v_1=\frac{v_s-v_1}{R_1}-\frac{v_1-v_2}{R_2},
\qquad
C_2\dot v_2=\frac{v_1-v_2}{R_2}.
\]
Substituindo os valores,
\[
\dot v_1=100u(t)-15v_1+5v_2,\qquad
\dot v_2=10v_1-10v_2.
\]
A matriz natural é
$A=\begin{bmatrix}-15&5\\10&-10\end{bmatrix}$, cuja equação característica é
$s^2+25s+100=0$, fornecendo $s_1=-5$ e $s_2=-20$. Em regime final,
$v_1(\infty)=v_2(\infty)=\SI{10}{\volt}$. Assim
$v_2=10+Ae^{-5t}+Be^{-20t}$. Das condições $v_2(0)=0$ e
$\dot v_2(0)=0$ resultam $A=-40/3$ e $B=10/3$.}
\end{questao}

\begin{questao}[4]{}
No circuito RLC em paralelo, a fonte de corrente vale $\SI{2}{\ampere}$ por muito
tempo e muda instantaneamente para $\SI{1}{\ampere}$ em $t=0$. Determine a
resposta completa $v(t)$, classifique o amortecimento e encontre o primeiro pico
negativo da tensão.

\circuitoqq{%
\begin{circuitikz}[scale=0.94,transform shape,line width=0.8pt,every node/.append style={font=\scriptsize}]
\draw (0,0) to[american current source] (0,3) -- (7.5,3);
\node[anchor=east,align=right] at (-0.25,1.5) {$I_s$\\$2\,\mathrm{A}\to1\,\mathrm{A}$};
\draw (2.3,3) to[R] (2.3,0);
\node[anchor=east] at (2.05,1.5) {$R=\SI{20}{\ohm}$};
\draw (4.9,3) to[C,v^>={$v$}] (4.9,0);
\node[anchor=east] at (4.65,1.5) {$C=\SI{500}{\micro\farad}$};
\draw (7.5,3) to[L,i>^={$i_L$}] (7.5,0) -- (0,0);
\node[anchor=east] at (7.25,1.5) {$L=\SI{0.2}{\henry}$};
\end{circuitikz}}{RLC em paralelo excitado por uma mudança de nível da fonte de corrente.}

\resposta{Subamortecida;
$v(t)\approx-23{,}09e^{-50t}\sin(86{,}60t)\ \si{\volt}$;
$t_p\approx\SI{12.1}{\milli\second}$;
$v_{\min}\approx-\SI{10.93}{\volt}$}
\resolucao{Antes da mudança, em regime CC, o indutor é curto e conduz toda a
corrente da fonte: $i_L(0^-)=\SI{2}{\ampere}$ e $v(0^-)=0$. Logo
$i_L(0^+)=\SI{2}{\ampere}$ e $v(0^+)=0$. Para $t>0$, o valor final é novamente
$v(\infty)=0$, com $i_L(\infty)=\SI{1}{\ampere}$. Os parâmetros são
\[
\alpha=\frac{1}{2RC}=\SI{50}{\per\second},\qquad
\omega_0=\frac{1}{\sqrt{LC}}=\SI{100}{\radian\per\second},
\]
portanto $\omega_d=\sqrt{100^2-50^2}\approx\SI{86.60}{\radian\per\second}$.
Aplicando KCL em $0^+$,
\[
1=\frac{0}{20}+C\dot v(0^+)+2
\Rightarrow
\dot v(0^+)=-\frac{1}{500\,\mathrm{\mu F}}
=-\SI{2000}{\volt\per\second}.
\]
Como $v(0)=0$, resulta
$v(t)=B e^{-50t}\sin(86{,}60t)$, com
$B=-2000/86{,}60\approx-23{,}09$. O primeiro extremo satisfaz
$\tan(\omega_dt)=\omega_d/\alpha=\sqrt3$, logo
$\omega_dt=\pi/3$ e $t_p\approx\SI{12.1}{\milli\second}$. Substituindo na
resposta, $v_{\min}\approx-\SI{10.93}{\volt}$.}
\end{questao}

\gabarito[Nível 4 --- projeto, diagnóstico e síntese]

